\section{Feature Computation and Preprocessing}

\subsection{Vegetation Index Derivation}
After acquiring multispectral surface reflectance imagery, the system computes vegetation indices that serve as quantitative indicators of canopy condition. The primary indices used are the Normalized Difference Vegetation Index (NDVI) and the Normalized Difference Moisture Index (NDMI), both of which are widely employed in forest health monitoring. NDVI captures greenness and photosynthetic activity, while NDMI reflects canopy moisture content—an important factor in early stress detection. These indices are computed directly from the red, near-infrared, and shortwave infrared bands of the Harmonized Landsat--Sentinel-2 (HLS) dataset. This approach aligns with prior studies showing that long-term NDVI and moisture-sensitive indices can reveal subtle stress signals before visible symptoms emerge (\cite{ndvi}).

\subsection{Forest Mask Application and Quality Filtering}
To ensure that index values represent true forest dynamics, the ESA WorldCover 10\,m land cover map is applied as a forest mask during preprocessing. This step removes cropland, grassland, water, and built-up surfaces, preventing misinterpretation of non-forest vegetation signals. The literature emphasizes the importance of such masking, noting that NDVI alone ``does not distinguish among vegetation types,'' and that restricting analysis to forest pixels improves interpretability and reduces contamination \cite{worldcover}. Additionally, cloud-affected observations are filtered using scene metadata and quality flags, ensuring that only valid, atmospherically corrected reflectance values contribute to the time series.

\subsection{Temporal Harmonization and Gap Handling}
Because satellite observations vary in frequency due to cloud cover, revisit cycles, and acquisition geometry, the preprocessing pipeline focuses on creating a consistent temporal structure for analysis. The Harmonized Landsat--Sentinel-2 (HLS) dataset already performs cross-sensor harmonization---standardizing reflectance values, aligning spectral bands, and resampling imagery onto a common grid---so no additional spectral harmonization is required during preprocessing. Instead, the pipeline aggregates observations into monthly composites or applies nearest-neighbor temporal alignment to reduce noise and ensure comparability across years. Missing observations caused by cloud contamination or data gaps are addressed through interpolation or omission depending on local data density. Establishing a stable, temporally consistent series is essential for constructing reliable baselines and aligns with the requirements of statistical time-series models such as BFAST \cite{bfast} and Adaptive EWMA \cite{forests-survey}, which depend on regular sampling to detect departures from expected seasonal behavior.

\subsection{Baseline Construction for Anomaly Detection}
For each forest pixel, a multi-year baseline is computed by aggregating monthly NDVI and NDMI values across the 10-year historical period. The baseline includes the mean and standard deviation for each month, forming a climatological signature of normal vegetation behavior. This structure mirrors the ``historical seasonal baseline'' approach used in near-real-time anomaly frameworks \cite{nrt-anomalies}, where deviations from expected monthly patterns are converted into standardized anomaly scores. The resulting baselines serve as the foundation for detecting statistically significant vegetation stress preceding disease outbreaks.
